\documentclass[../../main.tex]{subfiles}
\graphicspath{{../../image/}{../../figure/}}
\begin{document}

\section{热粘弹性}\label{section:ch02-热粘弹性}

现在我们考虑热弹性的一个推广，它涵盖了由\emph{率型粘性}诱导出\emph{内禀耗散}的材料。内能 $\varepsilon$、Piola--Kirchhoff 应力 $S$、温度 $\theta$ 和参考热通量向量 $Q$ 现在不仅可依赖于变形梯度 $F$、熵 $s$ 和温度梯度 $G$（如\eqref{align:ch02-2-5-1} 那样），还可依赖于变形梯度的时间变化率 $\dot{F}$：
\begin{align}
\begin{cases}
\varepsilon=\hat{\varepsilon}(F,\dot{F},s,G),\\
S=\hat{S}(F,\dot{F},s,G),\\
\theta=\hat{\theta}(F,\dot{F},s,G),\\
Q=\hat{Q}(F,\dot{F},s,G).
\end{cases}
\label{align:ch02-2-6-1}
\end{align}

如\cref{section:ch02-连续介质热力学平衡律} 所规定，任何平衡了质量、动量和能量的光滑热力学过程都必须恒等地满足耗散不等式\eqref{align:ch02-2-3-15}。将\eqref{align:ch02-2-6-1} 代入\eqref{align:ch02-2-3-15}，得
\begin{align}
\mathrm{tr}\,[(\rho_{0}\partial_{F}\hat{\varepsilon}-\hat{S})\dot{F}^{\top}]+\mathrm{tr}\,(\rho_{0}\partial_{\dot{F}}\hat{\varepsilon}\ddot{F}^{\top})+\rho_{0}(\partial_{s}\hat{\varepsilon}-\hat{\theta})\dot{s}+\rho_{0}\partial_{G}\hat{\varepsilon}\dot{G}-\hat{\theta}^{-1}\hat{Q}\cdot G\leqslant 0.
\label{align:ch02-2-6-2}
\end{align}

通过适当控制体力 $b$ 和热量供给 $r$，可以构造出平衡了质量、动量和能量的光滑过程，并在某点 $(x,t)$ 处对 $F,\dot{F},s,G,\ddot{F},\dot{s},\dot{G}$ 取任意预先指定的值，唯一约束是 $\det F>0$。因此不等式\eqref{align:ch02-2-6-2} 不可能恒成立，除非\eqref{align:ch02-2-6-1} 中的本构函数取如下特殊形式：
\begin{align}
\begin{cases}
\varepsilon=\hat{\varepsilon}(F,s),\\
S=\rho_{0}\partial_{F}\hat{\varepsilon}(F,s)+Z(F,\dot{F},s,G),\\
\theta=\partial_{s}\hat{\varepsilon}(F,s),\\
Q=\hat{Q}(F,\dot{F},s,G),
\end{cases}
\label{align:ch02-2-6-3}
\end{align}
\begin{align}
\mathrm{tr}\,[Z(F,\dot{F},s,G)\dot{F}^{\top}]+\frac{1}{\hat{\theta}(F,s)}\hat{Q}(F,\dot{F},s,G)\cdot G\geqslant 0.
\label{align:ch02-2-6-4}
\end{align}

将\eqref{align:ch02-2-6-3} 与 \eqref{align:ch02-2-5-3} 相比较，我们再次观察到：内能（它可仅依赖于变形梯度和熵）通过同样的 caloric 状态方程确定温度的本构方程。另一方面，应力的本构方程现在包含额外项 $Z$，它贡献粘性效应并诱导出由\eqref{align:ch02-2-6-4} 所体现的内禀耗散。

本构函数必须进一步约化，以遵守\cref{section:ch02-材料标架无差异} 中公设的材料标架无差异原理。特别地，标架无差异对内能施加与热弹性中相同的条件\eqref{align:ch02-2-5-5}$_{1}$，所得约化当然也相同：
\begin{align}
\hat{\varepsilon}(F,s)=\hat{\varepsilon}(U,s),
\label{align:ch02-2-6-5}
\end{align}
其中 $U$ 表示右拉伸张量\eqref{align:ch02-2-1-7}。此外，当\eqref{align:ch02-2-6-5} 成立时，通过\eqref{align:ch02-2-6-3}$_{3}$ 导出的温度本构方程以及应力本构方程中的项 $\rho_{0}\partial_{F}\hat{\varepsilon}(F,s)$ 将自动是标架无差异的。还需考察标架无差异对 $Z$ 和热通量的含义。由于分析最终将集中于热粘弹性流体，在此方便从 $S$、$Q$ 切换到 $T$、$q$；同时在状态变量清单 $(F,\dot{F},s,G)$ 中用 $L$（参见\eqref{align:ch02-2-1-8}）替换 $\dot{F}$，用 $g$（参见\eqref{align:ch02-2-3-17}）替换 $G$。于是我们写
\begin{align}
T=\rho\,\partial_{F}\hat{\varepsilon}(F,s)F^{\top}+\hat{Z}(F,L,s,g),
\label{align:ch02-2-6-6}
\end{align}
\begin{align}
q=\hat{q}(F,L,s,g).
\label{align:ch02-2-6-7}
\end{align}

回顾\eqref{align:ch02-2-4-3} 和 \eqref{align:ch02-2-4-5}，我们推出材料标架无差异原理要求
\begin{align}
\begin{cases}
\hat{Z}(OF,OLO^{\top}+\dot{O}O^{\top},s,Og)=O\hat{Z}(F,L,s,g)O^{\top}\\
\hat{q}(OF,OLO^{\top}+\dot{O}O^{\top},s,Og)=O\hat{q}(F,L,s,g),
\end{cases}
\label{align:ch02-2-6-8}
\end{align}
对任一正常正交矩阵 $O$ 成立。特别地，对任一固定状态 $(F,L,s,g)$ 及其自旋 $W$（参见\eqref{align:ch02-2-1-9}），可取 $O(t)=\exp(-tW)$，此时 $O(0)=I$，$\dot{O}(0)=-W$。由\eqref{align:ch02-2-6-8} 推得 $\hat{Z}$ 和 $\hat{q}$ 可仅通过其对称部分 $D$ 依赖于 $L$，因而\eqref{align:ch02-2-6-6} 和 \eqref{align:ch02-2-6-7} 可写为
\begin{align}
T=\rho\,\partial_{F}\hat{\varepsilon}(F,s)F^{\top}+\hat{Z}(F,D,s,g),
\label{align:ch02-2-6-9}
\end{align}
\begin{align}
q=\hat{q}(F,D,s,g),
\label{align:ch02-2-6-10}
\end{align}
其中 $\hat{Z}$ 和 $\hat{q}$ 满足
\begin{align}
\begin{cases}
\hat{Z}(OF,ODO^{\top},s,Og)=O\hat{Z}(F,D,s,g)O^{\top}\\
\hat{q}(OF,ODO^{\top},s,Og)=O\hat{q}(F,D,s,g),
\end{cases}
\label{align:ch02-2-6-11}
\end{align}
对所有正常正交矩阵 $O$ 成立。

为使角动量平衡律\eqref{align:ch02-2-3-9} 得到满足，$\hat{Z}$ 还必须是对称的：$\hat{Z}^{\top}=\hat{Z}$。注意此时耗散不等式\eqref{align:ch02-2-6-4} 可改写为
\begin{align}
\mathrm{tr}\,[\hat{Z}(F,D,s,g)D]+\frac{1}{\hat{\theta}(F,s)}\hat{q}(F,D,s,g)\cdot g\geqslant 0.
\label{align:ch02-2-6-12}
\end{align}

当介质被赋予材料对称性时，本构函数进一步约化。场在参考构型等容变化下的变换规则记录在\eqref{align:ch02-2-3-19}、\eqref{align:ch02-2-3-20}、\eqref{align:ch02-2-3-21} 和 \eqref{align:ch02-2-3-22} 中。如\cref{section:ch02-热弹性} 一样，我们在此引入材料的\emph{对称群} $\mathscr{G}$，即由那些使 $\varepsilon,T,\theta,q$ 的本构函数保持不变的参考构型等容变换 $\bar{x}$ 的 Jacobi 矩阵 $H$ 所构成的 $\mathrm{SL}(m)$ 子群。于是 $\mathscr{G}$ 是所有满足下式的 $H\in\mathrm{SL}(m)$ 之集：
\begin{align}
\begin{cases}
\hat{\varepsilon}(FH^{-1},s)=\hat{\varepsilon}(F,s),\\
\hat{Z}(FH^{-1},D,s,g)=\hat{Z}(F,D,s,g),\\
\hat{q}(FH^{-1},D,s,g)=\hat{q}(F,D,s,g).
\end{cases}
\label{align:ch02-2-6-13}
\end{align}

当 $\mathscr{G}\equiv\mathrm{SL}(m)$ 时，材料称为\emph{热粘弹性流体}。此时取 $H=(\det F)^{-1/m}F\in\mathrm{SL}(m)$，应用\eqref{align:ch02-2-6-13}，我们推出 $\hat{\varepsilon},\hat{Z},\hat{q}$ 可仅通过其行列式（或等价地，通过密度 $\rho$）依赖于 $F$。因此热粘弹性流体的本构方程约化为
\begin{align}
\begin{cases}
\varepsilon=\tilde{\varepsilon}(\rho,s),\\
T=-pI+\tilde{Z}(\rho,D,s,g),\\
p=\rho^{2}\partial_{\rho}\tilde{\varepsilon}(\rho,s),\quad\theta=\partial_{s}\tilde{\varepsilon}(\rho,s),\\
q=\tilde{q}(\rho,D,s,g).
\end{cases}
\label{align:ch02-2-6-14}
\end{align}

为满足标架无差异，$\tilde{Z}$ 和 $\tilde{q}$ 对任一正常正交矩阵 $O$ 还应满足条件
\begin{align}
\begin{cases}
\tilde{Z}(\rho,ODO^{\top},s,Og)=O\tilde{Z}(\rho,D,s,g)O^{\top},\\
\tilde{q}(\rho,ODO^{\top},s,Og)=O\tilde{q}(\rho,D,s,g),
\end{cases}
\label{align:ch02-2-6-15}
\end{align}
它们由\eqref{align:ch02-2-6-11} 推出。可以明确写出满足\eqref{align:ch02-2-6-15} 的最一般函数 $\tilde{Z}$ 和 $\tilde{q}$ 的形式。这里仅记录 $m=3$ 时与\eqref{align:ch02-2-6-15} 相容且关于 $(D,g)$ 线性的最一般本构关系：
\begin{align}
T=-p(\rho,s)I+\lambda(\rho,s)(\mathrm{tr}\,D)I+2\mu(\rho,s)D,
\label{align:ch02-2-6-16}
\end{align}
\begin{align}
q=\kappa(\rho,s)g,
\label{align:ch02-2-6-17}
\end{align}
它们标识了可压缩、导热的\emph{Newton 流体}。

Newton 流体的\emph{体积粘度} $\lambda+\tfrac{2}{3}\mu$、\emph{剪切粘度} $\mu$ 和\emph{热导率} $\kappa$ 受不等式\eqref{align:ch02-2-6-12} 约束，它在此约化为
\begin{align}
\lambda(\rho,s)(\mathrm{tr}\,D)^{2}+2\mu(\rho,s)\,\mathrm{tr}\,D^{2}+\frac{\kappa(\rho,s)}{\tilde{\theta}(\rho,s)}|g|^{2}\geqslant 0.
\label{align:ch02-2-6-18}
\end{align}
该不等式对任意 $D$ 和 $g$ 成立当且仅当
\begin{align}
\mu(\rho,s)\geqslant 0,\qquad 3\lambda(\rho,s)+2\mu(\rho,s)\geqslant 0,\qquad\kappa(\rho,s)\geqslant 0.
\label{align:ch02-2-6-19}
\end{align}
为产生实际耗散，$\mu$、$3\lambda+2\mu$ 和 $\kappa$ 中至少一个应严格为正。

\end{document}
